% IEEE journal article, from IEEEtran's bare_jrnl.tex (IEEEtran 1.8b).
% Use [journal,compsoc] for IEEE Computer Society journals and
% [journal,comsoc] for IEEE Communications Society journals.
\documentclass[journal]{IEEEtran}

\usepackage{cite}
\usepackage{amsmath,amssymb}
\usepackage{graphicx}

\begin{document}

\title{LATTICE_PROJECT_TITLE}

\author{First~Author,~\IEEEmembership{Member,~IEEE,}
        and~Second~Author% <- this % stops a space
\thanks{F. Author and S. Author are with the Department, Institution,
City, Country (e-mail: author@example.com).}% <- this % stops a space
\thanks{Manuscript received Month Day, Year; revised Month Day, Year.}}

% The running heads: journal name on the left, short author list and title
% on the right.
\markboth{IEEE Transactions on Journal Name}%
{Author \MakeLowercase{\textit{et al.}}: Short Title}

\maketitle

\begin{abstract}
Briefly summarize the problem, method, and main result.
\end{abstract}

\begin{IEEEkeywords}
first keyword, second keyword, third keyword
\end{IEEEkeywords}

\section{Introduction}
\IEEEPARstart{M}{otivate} the problem and state the article's main
contribution, building on prior work~\cite{shell2015ieeetran,kopka2004guide}.

\section{Related Work}
Position the contribution against the most relevant prior work.

\section{Method}
Describe the proposed method.

\section{Experiments}
Present the evaluation setup and results.

\section{Conclusion}
Summarize the findings and limitations.

\bibliographystyle{IEEEtran}
\bibliography{references}

\end{document}
